\subsection{阿廷模与诺特模}

定义 1. --- 设 A 为环。称 A-模 M 为阿廷模（相应地，诺特模），如果它满足下列等价条件：

\begin{enumerate}
\def\labelenumi{(\roman{enumi})}
\item
  M 的子模的任一非空集合按包含关系排序都有一个极小（相应地，极大）元。
\item
  M 的子模的任一递减（相应地，递增）序列都是平稳的。
\end{enumerate}

条件 (i) 与 (ii) 的等价性由《集合论》III, §6, No.~5, p.~190, 命题 6 得出。

A-模 M 是阿廷（相应地，诺特）的，当且仅当 M 看作位似环 \(A_{M}\) 上的模时是阿廷（相应地，诺特）的。

设 M 为阿廷（相应地，诺特）A-模。M 的子模的任一非空集合若按包含关系为左向定向（相应地，右向定向），则有一个最小元（相应地，最大元）（《集合论》III, §1, No.~10, p.~145, 命题 10）。

设 M 为阿廷（相应地，诺特）A-模，\((\mathrm{M}_{i})_{i\in\mathrm{I}}\) 为 M 的子模的一个族。族 \((\mathrm{M}_{i})_{i\in\mathrm{I}}\) 的有限子族的交（相应地，和）构成 M 的子模的一个非空左向（相应地，右向）定向集。因此，存在 I 的有限子集 J，使得 \(\bigcap_{i\in I}M_{i}=\bigcap_{i\in J}M_{i}\) （相应地，\(\sum_{i\in I}M_{i}=\sum_{i\in J}M_{i}\) ）。

例. --- 1) 体上的有限维向量空间既是阿廷的又是诺特的。

\begin{enumerate}
\def\labelenumi{\arabic{enumi})}
\setcounter{enumi}{1}
\tightlist
\item
  设 M 为 A-模。若存在 M 的非零子模的无穷族 \((\mathrm{M}_{i})_{i\in\mathrm{I}}\) ，其和是直和，则 M 既非阿廷亦非诺特：事实上，对 I 的子集的任一严格递减（相应地，严格递增）无穷序列 \((\mathrm{J}_{n})\) ，M 的子模的无穷序列 \((\sum_{i\in\mathrm{J}_{n}}\mathrm{M}_{i})\) 严格递减（相应地，严格递增）。特别地，体上的无穷维向量空间既非阿廷亦非诺特。
\end{enumerate}

*3) 我们将在后文看到，Z-模 Z 是诺特的但不是阿廷的（VIII, p.~5, 例 3）。*

\begin{enumerate}
\def\labelenumi{\arabic{enumi})}
\setcounter{enumi}{3}
\tightlist
\item
  设 \(p\) 为素数，\(\mathrm{M}_p\) 为 \(p\) -准素分量，即挠 \(\mathbf{Z}\) -模 \(\mathbf{Q} / \mathbf{Z}\) 的相应分量（VII, §2, No.~2, p.~7）。\(\mathrm{M}_p\) 的每个子模都等于 \(\mathrm{M}_p\) 或 \(p^{-n}\mathbf{Z} / \mathbf{Z}\) （\(n \in \mathbf{N}\) 为整数）（VII, §2, p.~54, 习题 3）。因此，\(\mathrm{M}_p\) 是阿廷的但非诺特的 \(\mathbf{Z}\) -模。
\end{enumerate}

命题 1. --- A-模 M 具有有限长度（II, §1, No.~10, p.~212），当且仅当它既是阿廷的又是诺特的。

设 M 具有有限长度 d。~于是 M 的子模的任一严格递增或严格递减序列至多有 \(d+1\) 项（I, §4, No.~7, p.~44）。因此，M 既是阿廷的又是诺特的。

反之，设 M 既是阿廷的又是诺特的。设 S 为 M 的具有有限长度的子模组成的集合。零子模是 S 的元素，且由于 M 是诺特的，S 有一个极大元 N。用反证法，设 \(M \neq N\) 。于是 M 的异于 N 且包含 N 的子模组成的集合有一个极小元 P，因为 M 是阿廷的。模 P/N 长度为 1，而由于 N 是有限长度的模，P 亦然（II, §1, No.~10, p.~212, 命题 16）。这与 N 的定义矛盾。

命题 2. --- A-模 M 是诺特的，当且仅当 M 的每个子模都是有限生成的。

先假设 M 的每个子模都是有限生成的。设 \((\mathrm{P}_{n})_{n\in\mathbf{N}}\) 为 M 的子模的一个递增序列，并设 P 为其并。它是 M 的一个子模。由假设，存在 M 的有限子集 F 生成模 P；设 \(n\in N\) 为使 \(F\subset P_{n}\) 的整数。于是有 \(P_{n}=P\) ，从而序列 \((\mathrm{P}_{n})_{n\in\mathbf{N}}\) 是平稳的。这就证明了模 M 是诺特的。

其逆是下面更精确的命题的推论。

引理 1. --- 设 M 为诺特 A-模，E 为 M 的子集。存在 E 的有限子集 F，它与 E 生成同一个子模。

事实上，由 VIII, p.~2，存在 E 的有限子集 F 使得 \(\sum_{x\in E}Ax=\sum_{x\in F}Ax\) 。

命题 3. --- 设 M 为 A-模，N 为 M 的子模。则 M 是阿廷（相应地，诺特）的，当且仅当 N 与 M/N 都是。

我们就阿廷模的情形给出证明；诺特模的情形是类似的。

设 M 是阿廷的。由于 N 的每个子模都是 M 的子模，模 N 是阿廷的。设 \((\mathrm{P}_{n})_{n\in\mathbf{N}}\) 为 M/N 的子模的一个递减序列。存在 M 的包含 N 的子模的递减序列 \((\mathrm{Q}_{n})_{n\in\mathbf{N}}\) ，使得 \(P_{n}=Q_{n}/N\) 对每个 \(n\in N\) 成立（I, §4, No.~6, p.~41, 定理 4）。由于 M 是阿廷的，序列 \((\mathrm{Q}_{n})\) 是平稳的，从而序列 \((\mathrm{P}_{n})\) 也是平稳的。因此，模 M/N 是阿廷的。

反之，设模 N 与 M/N 都是阿廷的，并考虑 M 的子模的一个递减序列 \((\mathrm{P}_{n})\) 。N 的子模序列 \(P_{n}^{\prime} = N \cap P_{n}\) 是平稳的。同样，M/N 的子模序列 \(\mathrm{P}_{n}^{\prime\prime} = (\mathrm{N} + \mathrm{P}_{n})/\mathrm{N}\) 也是平稳的。于是，存在整数 \(m \in N\) ，使得 \(P_{n}^{\prime} = P_{m}^{\prime}\) 与 \(P_{n}^{\prime\prime} = P_{m}^{\prime\prime}\) 对每个整数 \(n \geqslant m\) 成立。由下面的引理，序列 \((\mathrm{P}_{n})\) 因而是平稳的。

引理 2. --- 设 M 为 A-模，N、P、Q 为 M 的子模。假设有 \(P \subset Q\) ，\(N \cap P = N \cap Q\) 以及 \(N + P = N + Q\) 。则有 P = Q。

设 x 为 Q 的元素。它属于 \(N + P\) ；于是存在 P 的元素 y 使得 \(x - y \in N\) 。由于 Q 包含 P，差 x - y 属于 \(N \cap Q\) ，从而属于 P。因此，x 属于 P。

推论. --- 设 M 为 A-模，\((\mathrm{M}_{i})_{i\in\mathrm{I}}\) 为 M 的子模的有限族。

\begin{enumerate}
\def\labelenumi{\alph{enumi})}
\item
  若诸模 \(\mathrm{M}_i\) 都是阿廷（相应地，诺特）的，则它们的和 \(\sum_{i\in \mathrm{I}}\mathrm{M}_i\) 亦然。
\item
  若诸模 \(\mathrm{M} / \mathrm{M}_i\) 都是阿廷（相应地，诺特）的，则模 \(\mathrm{M} / \bigcap_{i\in \mathrm{I}}\mathrm{M}_i\) 亦然。
\end{enumerate}

用归纳法，只需处理 \(I = \{1, 2\}\) 的情形。模 \(M_{2}/(M_{1} \cap M_{2})\) 是 \(M_{2}\) 的商模，它同构于 \((M_{1} + M_{2})/M_{1}\) ，后者是 \(M/M_{1}\) 的子模（I, §4, No.~6, p.~41, 定理 4）。

在 a) 中，设 \(\mathrm{M}_1\) 与 \((\mathrm{M}_1 + \mathrm{M}_2) / \mathrm{M}_1\) 是阿廷（相应地，诺特）的；于是 \(\mathrm{M}_1 + \mathrm{M}_2\) 亦然（命题 3）。

在 b) 中，设 \(M/M_{2}\) 与 \(M_{2}/(M_{1} \cap M_{2})\) 是阿廷（相应地，诺特）的；于是 \(M/(M_{1} \cap M_{2})\) 亦然（loc. cit.）。

例 5. --- 设 \((\mathrm{M}_{i})_{i\in\mathrm{I}}\) 为 A-模的有限族。若诸模 \(M_{i}\) 都是阿廷（相应地，诺特）的，则它们的直和 \(\bigoplus_{i\in I} M_{i}\) 亦然。

注. --- 本小节的定义和结果可推广到带算子的任意交换群，只需把陈述中的子模换为稳定子群。

